\section{Fourier rigidity and the central excess}
\label{sec:s5-crystal-rigidity}

The argument here uses the actual signed divisor and its ordinary
distributional Fourier transform.  It does not assume that the
divisor is positive, that zeros are simple, or that the degree-ten
symmetric-square function is automorphic.  The analytic input is
the positive-definite atomic coefficient function constructed in
MA-FQ-001.  The finite-coset reduction is MA-FQ-002.

\paragraph{MA-FQ-003 (integer atoms and a prime-frequency obstruction).}
Let $a:\mathbb R\to\mathbb Z$ be positive definite when $\mathbb R$
is regarded as a discrete additive group.  Suppose its nonzero
support is locally finite in the ordinary real topology, and
the atomic measure
\[
                 \nu=\sum_{t\in\mathbb R}a(t)\delta_t
\tag{MA-FQ.20}
\]
is tempered.  Suppose also that its ordinary Fourier transform,
with convention
$\widehat\nu(x)=\int e^{itx}\,\nu(dt)$, is a purely atomic
locally finite complex measure.  Then there are finitely many
numbers $h_1,\ldots,h_N>0$ such that
\[
                 \operatorname{supp}\widehat\nu
                        \subseteq\bigcup_{j=1}^{N}h_j\mathbb Z.
\tag{MA-FQ.21}
\]
The empty union is allowed when $\nu=0$.
In particular, for any fixed $c>0$, only finitely many distinct
rational primes $p$ can have
\[
                         \widehat\nu(\{c\log p\})\ne0.
\tag{MA-FQ.22}
\]

\begin{proof}
The two-by-two positive-definite inequality gives
$|a(t)|\le a(0)$ for every $t$.  If $a(0)=0$, then $a=0$
and the conclusions follow.  Otherwise put $k=a(0)\in\mathbb N$.
The function has finite integer range contained in
$\{-k,-k+1,\ldots,k\}$.
By Bochner's theorem on the locally compact discrete abelian
group $\mathbb R_d$, it belongs to $B(\mathbb R_d)$, with
$\|a\|_{B(\mathbb R_d)}=k$.  Its representing measure lives on
the compact Pontryagin dual of $\mathbb R_d$; it is not being
identified with the ordinary Fourier transform in (MA-FQ.21).

For every nonzero value $r$ in the range, polynomial
interpolation on this finite range gives a polynomial $P_r$
with $P_r(r)=1$ and $P_r(q)=0$ at every other range value.
Thus $P_r(a)=1_{\{a=r\}}\in B(\mathbb R_d)$, since the
Fourier--Stieltjes algebra is closed under multiplication
and contains constants.  Cohen's idempotent theorem places
each such level set in the coset ring of $\mathbb R_d$.
These level sets are locally finite.  MA-FQ-002 therefore
writes their indicators, and hence $a$, as a finite integer
linear combination of indicators of ordinary discrete
cosets in $\mathbb R$.

Every nonzero discrete subgroup of $\mathbb R$ is of the
form $\alpha\mathbb Z$ with $\alpha>0$; the zero subgroup
gives a singleton.  Combining repeated terms, we obtain
an equality of tempered measures
\[
 \nu=\sum_{j=1}^{J}c_j\sum_{n\in\mathbb Z}
                         \delta_{\alpha_j n+\beta_j}
                  +\sum_{\ell=1}^{L}d_\ell\delta_{\tau_\ell},
 \qquad c_j,d_\ell\in\mathbb Z,\quad \alpha_j>0 .
\tag{MA-FQ.23}
\]
No positivity of the individual coefficients in this
representation is asserted or needed.  In particular,
the two Fourier transforms used in the preceding paragraph
and the next paragraph are kept distinct.

The ordinary Poisson summation identity for a lattice gives
\[
 \widehat{\sum_{n\in\mathbb Z}\delta_{\alpha n+\beta}}
    =\frac{2\pi}{\alpha}\sum_{m\in\mathbb Z}
        e^{2\pi i m\beta/\alpha}\delta_{2\pi m/\alpha}.
\tag{MA-FQ.24}
\]
This is an identity of tempered distributions, obtained by
pairing with a Schwartz test and applying the classical
one-dimensional Poisson formula.  A singleton transforms
to the ordinary continuous function $e^{i\tau x}$, interpreted
as its density against $dx$.  Thus (MA-FQ.23) transforms to
a measure supported on finitely many unshifted frequency
lattices, plus the absolutely continuous measure with density
\[
                         \sum_{\ell=1}^{L}d_\ell e^{i\tau_\ell x}.
\]
By hypothesis the complete transform is purely atomic.
Uniqueness of the local Lebesgue decomposition forces this
continuous density to vanish almost everywhere, hence
everywhere by continuity.  The remaining support is
contained in the finite union of lattices with
$h_j=2\pi/\alpha_j$, proving (MA-FQ.21).
Any cancellation among the lattice contributions can only
shrink their union of supports.

For a fixed $h>0$, suppose $c\log p=nh$ and
$c\log q=mh$ for distinct primes $p,q$ and positive
integers $n,m$.  Then $m\log p=n\log q$, so
$p^m=q^n$, contradicting unique factorization in
$\mathbb Z$.  Each lattice in (MA-FQ.21) therefore
contains $c\log p$ for at most one distinct prime.
This proves (MA-FQ.22).
\end{proof}

\paragraph{MA-FQ-004 (central equality in the stipulated family under
$(\mathcal R_B)$).}
Retain the exact $S_5$ family and the hypothesis
$(\mathcal R_B)$ of MA-S5-001.  Then
\[
                              \boxed{\beta=b+n.}
\tag{MA-FQ.25}
\]
Equivalently, the exceptional set $\Omega$ is finite and
the two rank-five comparison factors match at every place,
in the precise sense of MT-S508 and MA-S5-001.

\begin{proof}
MA-S5-001 first gives $k=\beta-b-n\ge0$.
Suppose for contradiction that $k>0$.
Use the actual difference datum
$\mathcal F_R=\mathcal F_{\varpi}-\mathcal F_W$ from
MA-LD-001 and its signed completed divisor
\[
              \mu=\sum_{t\in\mathbb R}m_R(t)\delta_t.
\]
Its support is on the critical line, written here in
height coordinates, by MA-LD-001.  That result uses
$(\mathcal R_B)$ and the positive-excess exceptional-prime
asymptotic, not the input $M$.
In particular $m_R(0)=-k$.

MA-FQ-001 proves that
$a(t)=-m_R(t)$ is an integer-valued positive-definite
function on $\mathbb R_d$.  It has locally finite support,
and $-\mu$ is tempered.  The same proof verifies that its
ordinary Fourier transform is a locally finite purely
atomic measure:
\[
                       \widehat{-\mu}=M_{\rm good}+B.
\tag{MA-FQ.26}
\]
Here the possible conductor atom and every bad-prime
correction are contained in the finite-total-variation
measure $B$.  The good-prime part has, for each
$p\in\Omega$, the nonzero atom
\[
              M_{\rm good}(\{\log p\})
                        =4(\log p)p^{-1/2}>0.
\tag{MA-FQ.27}
\]
The set $\Omega$ contains only good primes.  No power
of a different rational prime has norm $p$, so neither
another prime-power term nor a bad-prime correction
cancels this atom.  The conductor atom is at zero.
Consequently $\widehat{-\mu}(\{\log p\})\ne0$ for every
$p\in\Omega$.

Apply MA-FQ-003 to $\nu=-\mu$ with $c=1$.
It follows that $\Omega$ is finite.
MA-S5-001, the fully checked form of MT-S508,
then gives $k=0$, contradicting $k>0$.
Thus $k=0$ and (MA-FQ.25) follows.  The equivalent
all-place comparisons are the conclusion of that same
checked source criterion.
\end{proof}

\paragraph{Scope of the conclusion and the other inputs.}
The hypothesis $(\mathcal R_B)$ remains explicit.
This theorem does not prove it for every extension in the
family.  It excludes an additional central excess under
exactly that hypothesis, so it closes the central branch
that remained in MT-S508.  The all-place comparison gives
the Artin automorphy and holomorphy consequences recorded
in MA-S5-001.

GRH for all irreducible factors still requires the
appropriate Hecke zero-location input, sufficiently $M$.
Since $B$ is a Hecke function, $M$ also supplies
$(\mathcal R_B)$ in this family.  The Dirichlet theorem
$A$ alone does not replace that general Hecke premise.
No field with $k>0$ is constructed, and no unrestricted
$S_5$ Artin conjecture is claimed.

MA-LD-001--004 and the formal order assignment MA-S5-005
remain valid conditional deductions.  Their positive-excess
antecedent is now excluded by MA-FQ-004.  The formal
assignment was never an arithmetic realization; the
Fourier rigidity condition is the additional structure
that the central-order inequalities did not encode.
Similarly, the separate sufficient functoriality criterion
MA-S5-007 remains correct but is no longer needed to
eliminate this central branch.

For Artin motives, the conclusion identifies the actual
rank-five motive's $L$-function with its automorphic
comparison under the stipulated family hypotheses.
The finite-monodromy interpretation is the one in
SGA~1, Expos\'e~V, \S7 and Deligne,
\emph{Valeurs de fonctions $L$}, Definition~6.1,
\S\S6.2--6.3.  The proof uses harmonic analysis of the
signed divisor and prime frequencies in addition to
representation-ring identities.  It does not construct
a positive faithful Hilbert action on all zero jets,
and therefore does not evade or contradict MT-M20/M23.

\paragraph{Primary sources and proof-check boundary.}
The only deep harmonic-analysis input to MA-FQ-003 is
the classical idempotent theorem and its discrete-coset
consequence, with the precise versions checked in
MA-FQ-002.  An exact modern primary statement is
Sanders (2020), \S12, Theorems~12.1 and~12.3,
\emph{Bounds in Cohen's idempotent theorem},
\emph{Journal of Fourier Analysis and Applications, 26},
Article~25,
\url{https://doi.org/10.1007/s00041-020-09732-y}.
The passage from ordinary discrete cosets to lattice
combs and then to prime-frequency finiteness is proved
in full above.

The discrete-coset theorem is also stated and proved in
Kolountzakis (2000), \S3.2, Lemma~2 and Theorem~3,
\emph{On the structure of multiple translational tilings
by polygonal regions},
\emph{Discrete \& Computational Geometry, 23}, 537--553;
author version \url{https://arxiv.org/abs/math/9904065},
pp.~7--8.  The present argument uses the theorem in one
real dimension, with the ordinary topology used for
local finiteness and the discrete topology used for
the Fourier--Stieltjes algebra.
